% Manuscript-ready fragment. Bibliography keys: Evlampiev2009, Gainov2010.
% Archival evidence and subsequent author recollection are explicitly separated.
\subsection{Low-frequency NQR evidence}

The thesis explicitly reports copper NQR signals at 4.2~K in the
1.5--2.2~MHz interval \citep[p.~22]{Evlampiev2009}.
Its Fig.~8, captioned for a sample of nominal composition
Cu$_{1.1}$S, shows both a frequency-stepped spin-echo spectrum and an
FT-based reconstruction \citep[p.~23]{Evlampiev2009}.
The nominal composition describes a mixed synthetic sample, rather than
the stoichiometry of an individual anilite phase.
The thesis states that broadening and overlap prevented separation of
the component then called geerite from covellite and prevented a count
of the corresponding inequivalent Cu sites.
Its discussion already tentatively associated the low-frequency response
with tetrahedrally coordinated copper \citep[p.~33]{Evlampiev2009}.
The spectrum was therefore not assigned in its entirety to an unknown
extraneous impurity.

A related review retained the tentative entry ``Site~?'' and
``Below 2~MHz?'' for the material then called geerite
\citep[Table~I, p.~332]{Gainov2010}.
The same source gives the covellite low-frequency pair at 1.87 and
1.73~MHz for $^{63}$Cu and $^{65}$Cu, respectively, at 4.2~K.
Consequently, the archived low-frequency intensity cannot be attributed
entirely to anilite.

The experimenter additionally reports that the low-frequency
instrumentation and resonant circuits had not been sufficiently optimized in
this range for a detailed investigation (author communication,
27 November 2026). This recollection supplements the archive:
the thesis lists a nominal 1--32~MHz low-frequency spectrometer, but its
explicit reason for not resolving the low-frequency components is their
width and overlap.
The earlier publication independently documents an effective lower
frequency limit of about 1.5~MHz and reports that the covellite
low-frequency response could be followed only up to 30~K
\citep{Gainov2009}. This published operating limit supports the existence
of instrumental constraints, while the recollection about circuit
optimization remains a separate source of information.

Under the revised phase identification, the small model EFG at Cu3
provides a site-specific hypothesis for a possible anilite contribution
within this previously observed interval. It does not establish the
frequency of Cu3, resolve its isotope pair, or quantify its contribution
to the mixed spectrum. The observation supports a targeted assignment
test while retaining covellite as a known overlapping response.
